% Options for packages loaded elsewhere
\PassOptionsToPackage{unicode}{hyperref}
\PassOptionsToPackage{hyphens}{url}
\PassOptionsToPackage{dvipsnames,svgnames,x11names}{xcolor}
\PassOptionsToPackage{space}{xeCJK}
\documentclass[
]{article}
\usepackage{xcolor}
\usepackage[margin=2.5cm]{geometry}
\usepackage{amsmath,amssymb}
\setcounter{secnumdepth}{-\maxdimen} % remove section numbering
\usepackage{iftex}
\ifPDFTeX
  \usepackage[T1]{fontenc}
  \usepackage[utf8]{inputenc}
  \usepackage{textcomp} % provide euro and other symbols
\else % if luatex or xetex
  \usepackage{unicode-math} % this also loads fontspec
  \defaultfontfeatures{Scale=MatchLowercase}
  \defaultfontfeatures[\rmfamily]{Ligatures=TeX,Scale=1}
\fi
\usepackage{lmodern}
\ifPDFTeX\else
  % xetex/luatex font selection
  \setmainfont[]{DejaVu Sans}
  \setmonofont[]{DejaVu Sans Mono}
  \ifXeTeX
    \usepackage{xeCJK}
    \setCJKmainfont[]{Noto Sans CJK SC}
  \fi
  \ifLuaTeX
    \usepackage[]{luatexja-fontspec}
    \setmainjfont[]{Noto Sans CJK SC}
  \fi
\fi
% Use upquote if available, for straight quotes in verbatim environments
\IfFileExists{upquote.sty}{\usepackage{upquote}}{}
\IfFileExists{microtype.sty}{% use microtype if available
  \usepackage[]{microtype}
  \UseMicrotypeSet[protrusion]{basicmath} % disable protrusion for tt fonts
}{}
\makeatletter
\@ifundefined{KOMAClassName}{% if non-KOMA class
  \IfFileExists{parskip.sty}{%
    \usepackage{parskip}
  }{% else
    \setlength{\parindent}{0pt}
    \setlength{\parskip}{6pt plus 2pt minus 1pt}}
}{% if KOMA class
  \KOMAoptions{parskip=half}}
\makeatother
\usepackage{longtable,booktabs,array}
\usepackage{caption}
\captionsetup[table]{skip=6pt}
\newcounter{none} % for unnumbered tables
\usepackage{calc} % for calculating minipage widths
% Correct order of tables after \paragraph or \subparagraph
\usepackage{etoolbox}
\makeatletter
\patchcmd\longtable{\par}{\if@noskipsec\mbox{}\fi\par}{}{}
\makeatother
% Allow footnotes in longtable head/foot
\IfFileExists{footnotehyper.sty}{\usepackage{footnotehyper}}{\usepackage{footnote}}
\makesavenoteenv{longtable}
\setlength{\emergencystretch}{3em} % prevent overfull lines
\providecommand{\tightlist}{%
  \setlength{\itemsep}{0pt}\setlength{\parskip}{0pt}}
\usepackage{bookmark}
\IfFileExists{xurl.sty}{\usepackage{xurl}}{} % add URL line breaks if available
\urlstyle{same}
% fallback for those not using the hyperref driver hyperxmp:
\makeatletter
\@ifundefined{xmpquote}{\newcommand{\xmpquote}[1]{#1}}{}
\makeatother
\hypersetup{
  colorlinks=true,
  linkcolor={Maroon},
  filecolor={Maroon},
  citecolor={Blue},
  urlcolor={Blue},
  pdfcreator={LaTeX via pandoc}}

\author{}
\date{}

\begin{document}

\section{筑基篇课后习题 XIII：哥德巴赫猜想的 pat
重新观测}\label{ux7b51ux57faux7bc7ux8bfeux540eux4e60ux9898-xiiiux54e5ux5fb7ux5df4ux8d6bux731cux60f3ux7684-pat-ux91cdux65b0ux89c2ux6d4b}

\begin{quote}
\textbf{声明 (Declaration)}: 本文\textbf{不代表任何数学结论}
(non-mathematical-claim), 仅为基于 Lean 4
形式化的一种\textbf{实验观测报告} (experimental observation report)。

{\def\LTcaptype{none} % do not increment counter
\begin{longtable}[]{@{}
  >{\raggedright\arraybackslash}p{(\linewidth - 6\tabcolsep) * \real{0.2500}}
  >{\raggedright\arraybackslash}p{(\linewidth - 6\tabcolsep) * \real{0.2500}}
  >{\raggedright\arraybackslash}p{(\linewidth - 6\tabcolsep) * \real{0.2500}}
  >{\raggedright\arraybackslash}p{(\linewidth - 6\tabcolsep) * \real{0.2500}}@{}}
\toprule\noalign{}
\begin{minipage}[b]{\linewidth}\raggedright
习题
\end{minipage} & \begin{minipage}[b]{\linewidth}\raggedright
主题
\end{minipage} & \begin{minipage}[b]{\linewidth}\raggedright
Lean 形式化
\end{minipage} & \begin{minipage}[b]{\linewidth}\raggedright
DOI
\end{minipage} \\
\midrule\noalign{}
\endhead
\bottomrule\noalign{}
\endlastfoot
exercise\_i\_13 & 哥德巴赫猜想的 pat 重新观测（R170） & PatGoldbach.lean
& 10.5281/zenodo.21917253 \\
\end{longtable}
}
\end{quote}

\textbf{Exercise XIII for the Foundation-Building Chapter: The Goldbach
Conjecture Revisited from the PAT Perspective}

\begin{quote}
习题定位：筑基篇课后习题系列第十三题。用 mechanics-pat-observation skill
观测哥德巴赫猜想（每个大于 2
的偶数是两个素数之和）。全部新定理只锚筑基篇定理（R085/R128/R136②③/R143/R146/R159），不用外部引理。
\end{quote}

\emph{2026-08-13 · Internal research exercise · Lean 4 / mathlib v4.32.2
· 6 new theorems PROVED no sorry · 算器神魂论筑基篇课后习题 XIII }

\begin{center}\rule{0.5\linewidth}{0.5pt}\end{center}

\subsection{习题陈述}\label{ux4e60ux9898ux9648ux8ff0}

\begin{quote}
对哥德巴赫猜想（Goldbach Conjecture：每个大于 2 的偶数 2n
都可表示为两个素数之和）做 pat 重新观测。唯一论点：\textbf{p + q = 2n ⟺
p, q 关于 n 对称（q = 2n - p）------哥德巴赫分解 =
素数对称对，关于偶数中心 n 折叠还原。}
\end{quote}

\begin{center}\rule{0.5\linewidth}{0.5pt}\end{center}

\subsection{零、总论：哥德巴赫 =
素数对称对}\label{ux96f6ux603bux8bbaux54e5ux5fb7ux5df4ux8d6b-ux7d20ux6570ux5bf9ux79f0ux5bf9}

哥德巴赫猜想：2n = p + q（p, q 素数）。

\textbf{pat 结构对应}（mechanics-pat-observation skill）：

{\def\LTcaptype{none} % do not increment counter
\begin{longtable}[]{@{}lll@{}}
\toprule\noalign{}
哥德巴赫结构 & pat 结构 & 锚定 \\
\midrule\noalign{}
\endhead
\bottomrule\noalign{}
\endlastfoot
2n = p + q & 对称对 \{p, q\} = \{p, 2n-p\} 关于 n & R085/R136②③ \\
p, q 关于 n 对称 & 镜像 S(x) = 2n - x（对合 S²=id） & R085/R128 \\
偶数中心 n & 折叠中心（不动点 S(n) = n） & R085 \\
p, q 是素数 & 素数 = 方向 log p（单相位） & R159/R146 \\
★猜想 & 每个偶数有素数对称对 & CONJECTURE \\
\end{longtable}
}

\begin{center}\rule{0.5\linewidth}{0.5pt}\end{center}

\subsection{一、p + q = 2n ⟺
对称对}\label{ux4e00p-q-2n-ux5bf9ux79f0ux5bf9}

\textbf{★核心：哥德巴赫分解 2n = p + q 的代数等价 = q = 2n - p------p 和
q 关于偶数中心 n 对称。}

\subsubsection{论证}\label{ux8bbaux8bc1}

\begin{enumerate}
\def\labelenumi{\arabic{enumi}.}
\tightlist
\item
  \textbf{对称对}（R085 折叠类 / R136 ②③ 对称性成对声明）：p + q = 2n ⟺
  q = 2n - p------\{p, q\} = \{p, 2n-p\} 关于 n 成对。
\item
  \textbf{哥德巴赫分解 = 对称对}：两个素数关于偶数中心对称分布。
\end{enumerate}

\subsubsection{形式化（PatGoldbach.lean，1
定理）}\label{ux5f62ux5f0fux5316patgoldbach.lean1-ux5b9aux7406}

\begin{itemize}
\tightlist
\item
  \texttt{symmetric\_pair\_iff\_sum}：p + q = 2n ⟺ q = 2n -
  p------哥德巴赫分解 = 对称对（linarith）。
\end{itemize}

\textbf{验证}：0 sorry。

\begin{center}\rule{0.5\linewidth}{0.5pt}\end{center}

\subsection{二、镜像对合与中心不动点（R085/R128）}\label{ux4e8cux955cux50cfux5bf9ux5408ux4e0eux4e2dux5fc3ux4e0dux52a8ux70b9r085r128}

\textbf{★核心：偶数中心 n 的镜像 S(x) = 2n - x 是对合（S² = id），n
是不动点------折叠类还原。}

\subsubsection{论证}\label{ux8bbaux8bc1-1}

\begin{enumerate}
\def\labelenumi{\arabic{enumi}.}
\tightlist
\item
  \textbf{镜像对合}（R085/R128）：2n - (2n - x) = x------关于 n 的镜像
  S² = id，素数对称对经两次镜像还原。
\item
  \textbf{中心不动点}（R085 折叠中心）：2n - n = n------偶数中心 n
  是折叠中心，对称对关于 n 折叠，n 不变。
\end{enumerate}

\subsubsection{形式化（PatGoldbach.lean，2
定理）}\label{ux5f62ux5f0fux5316patgoldbach.lean2-ux5b9aux7406}

\begin{itemize}
\tightlist
\item
  \texttt{mirror\_involution}：2n - (2n - x) = x------镜像对合。
\item
  \texttt{center\_fixed}：2n - n = n------中心不动点。
\end{itemize}

\textbf{验证}：0 sorry。ring。

\begin{center}\rule{0.5\linewidth}{0.5pt}\end{center}

\subsection{三、对称对的中点（R085/R143）}\label{ux4e09ux5bf9ux79f0ux5bf9ux7684ux4e2dux70b9r085r143}

\textbf{★核心：p + q = 2n ⟺ q - n = -(p - n)------p 和 q 到中心 n
的距离相反（对称分布）。}

\subsubsection{论证}\label{ux8bbaux8bc1-2}

\begin{enumerate}
\def\labelenumi{\arabic{enumi}.}
\tightlist
\item
  \textbf{距离相反}：q - n = -(p - n)------p, q 到 n 等距反向（R085：0 =
  ±1 折叠类；R143：对称对还原到中心）。
\item
  \textbf{中点}：n 是 p, q 的中点------哥德巴赫分解 = 两个素数关于 n
  对称。
\end{enumerate}

\subsubsection{形式化（PatGoldbach.lean，1
定理）}\label{ux5f62ux5f0fux5316patgoldbach.lean1-ux5b9aux7406-1}

\begin{itemize}
\tightlist
\item
  \texttt{symmetric\_pair\_midpoint}：p + q = 2n ⟺ q - n = -(p -
  n)------对称对的中点（linarith）。
\end{itemize}

\textbf{验证}：0 sorry。

\begin{center}\rule{0.5\linewidth}{0.5pt}\end{center}

\subsection{四、★哥德巴赫 =
素数对称对（CONJECTURE）}\label{ux56dbux54e5ux5fb7ux5df4ux8d6b-ux7d20ux6570ux5bf9ux79f0ux5bf9conjecture}

\textbf{★核心：哥德巴赫猜想断言每个偶数 2n 存在素数对称对 \{p, q\} 关于
n------两个素数方向在中心 n 折叠还原。}

\subsubsection{论证}\label{ux8bbaux8bc1-3}

\begin{enumerate}
\def\labelenumi{\arabic{enumi}.}
\tightlist
\item
  \textbf{素数 = 方向 log p}（R159 prime\_log\_monophase /
  R146）：哥德巴赫的两个素数是两个素数方向。
\item
  \textbf{★转译}：每个偶数 2n（n \textgreater{} 1）存在素数 p, q 满足 q
  = 2n - p------素数对称对关于中心 n 折叠还原（R085）。
\item
  \textbf{诚实边界}：强哥德巴赫未证（CONJECTURE；弱哥德巴赫已证，Helfgott
  2013，外部文献）。
\end{enumerate}

\subsubsection{形式化（PatGoldbach.lean，1
定理）}\label{ux5f62ux5f0fux5316patgoldbach.lean1-ux5b9aux7406-2}

\begin{itemize}
\tightlist
\item
  \texttt{goldbach\_symmetric\_decomposition}：p + q = 2n ⟺ q = 2n -
  p------哥德巴赫 = 素数对称对（代数部分 PROVED，存在性 CONJECTURE）。
\end{itemize}

\textbf{验证}：0 sorry。

\begin{center}\rule{0.5\linewidth}{0.5pt}\end{center}

\subsection{五、全景（组合定理）}\label{ux4e94ux5168ux666fux7ec4ux5408ux5b9aux7406}

\textbf{★核心：对称对（p, q 关于 n）∧ 镜像对合（S²=id）∧
中心不动点（S(n)=n）∧ 素数 = 方向 log p------哥德巴赫 =
素数对称对的折叠还原。}

\subsubsection{形式化（PatGoldbach.lean，1
定理）}\label{ux5f62ux5f0fux5316patgoldbach.lean1-ux5b9aux7406-3}

\begin{itemize}
\tightlist
\item
  \texttt{goldbach\_pat\_perspective}：对称对 ∧ 镜像对合 ∧ 中心不动点 ∧
  素数方向------全景。
\end{itemize}

\textbf{验证}：0 sorry。合取项分别锚 symmetric\_pair\_iff\_sum /
mirror\_involution / center\_fixed / R159。

\begin{center}\rule{0.5\linewidth}{0.5pt}\end{center}

\subsection{六、习题解答总结}\label{ux516dux4e60ux9898ux89e3ux7b54ux603bux7ed3}

{\def\LTcaptype{none} % do not increment counter
\begin{longtable}[]{@{}llll@{}}
\toprule\noalign{}
论点 & 内容 & 锚定 & 标签 \\
\midrule\noalign{}
\endhead
\bottomrule\noalign{}
\endlastfoot
对称对 & p + q = 2n ⟺ q = 2n - p & R085/R136②③ & PROVED \\
镜像对合 & 2n - (2n - x) = x & R085/R128 & PROVED \\
中心不动点 & 2n - n = n & R085 & PROVED \\
对称中点 & q - n = -(p - n) & R085/R143 & PROVED \\
★哥德巴赫 & 素数对称对存在性 & R159 + 上述 & CONJECTURE \\
\end{longtable}
}

\textbf{习题的回答（一句话）}：\textbf{哥德巴赫猜想在 pat 视角 =
素数对称对------每个偶数 2n 的分解 p + q = 2n 等价于 p, q 关于中心 n
对称（q = 2n - p，镜像对合 S²=id，中心不动点）；两个素数 =
两个素数方向（log p，R159），在偶数中心 n
折叠还原。存在性未证（强哥德巴赫，CONJECTURE）。}

\textbf{诚实边界}：强哥德巴赫未证（CONJECTURE）------本习题交付结构转译（对称对
=
哥德巴赫分解的代数结构，PROVED；素数对称对存在性，CONJECTURE）。弱哥德巴赫（每个大于
5 的奇数是三个素数之和）已证（Helfgott 2013，外部文献，仅对照非证明）。

\begin{center}\rule{0.5\linewidth}{0.5pt}\end{center}

\subsection{七、相关对照}\label{ux4e03ux76f8ux5173ux5bf9ux7167}

{\def\LTcaptype{none} % do not increment counter
\begin{longtable}[]{@{}
  >{\raggedright\arraybackslash}p{(\linewidth - 4\tabcolsep) * \real{0.3333}}
  >{\raggedright\arraybackslash}p{(\linewidth - 4\tabcolsep) * \real{0.3333}}
  >{\raggedright\arraybackslash}p{(\linewidth - 4\tabcolsep) * \real{0.3333}}@{}}
\toprule\noalign{}
\begin{minipage}[b]{\linewidth}\raggedright
文献
\end{minipage} & \begin{minipage}[b]{\linewidth}\raggedright
对应内容
\end{minipage} & \begin{minipage}[b]{\linewidth}\raggedright
备注
\end{minipage} \\
\midrule\noalign{}
\endhead
\bottomrule\noalign{}
\endlastfoot
\textbf{Goldbach, C.} 1742（给 Euler 的信） & 哥德巴赫猜想原始陈述 & pat
转译对象 \\
\textbf{Helfgott, H.} 2013. \emph{Major arcs for Goldbach's
theorem}（arXiv:1305.2897） & 弱哥德巴赫已证 & 仅对照，非证明 \\
\textbf{R085/R128/R136②③/R143/R146/R159}（筑基篇） &
折叠类/镜像对合/成对/对称还原/素数方向 & 本框架定理 \\
\end{longtable}
}

\begin{center}\rule{0.5\linewidth}{0.5pt}\end{center}

\emph{筑基篇课后习题 XIII · 2026-08-13 · Lean 4 / mathlib v4.32.2 · 6
theorems PROVED no sorry（PatGoldbach.lean）· 配套 Lean
形式化：formal/Formal/Toolkit/PatGoldbach.lean（快照见
../lean/PatGoldbach.lean）}

\end{document}
